% Responsable: Integrante 2 (tarea 2.4 del plan)
\section{Diseño del compensador}\label{sec:diseno}
\subsection{Estructura y ley de control}
El controlador PI se escribe como
\begin{equation}
C(s)=K_c\,\frac{s+z}{s}=K_c\left(1+\frac{1}{T_is}\right),\qquad z=\frac{1}{T_i},\label{eq:PI}
\end{equation}
y, en el dominio del tiempo, la corriente aplicada a la celda es
\begin{equation}
\begin{split}
I(t)&=\mathrm{sat}_{[0,\,2\,\mathrm A]}\Big\{I_0+K_c\,e(t)+K_i\!\int_0^t e(\tau)\,d\tau\Big\},\\
e(t)&=T_b(t)-T_{\rm ref},
\end{split}\label{eq:ley}
\end{equation}
con $K_i=K_c/T_i$. El término $I_0$ es la corriente del punto de operación y actúa como valor inicial de la salida. La Figura~\ref{fig:lazoPI} muestra cómo se implementa el lazo en Simulink.

\begin{figure}[!tbp]\centering
\begin{tikzpicture}[>=Latex,font=\footnotesize,
  bloque/.style={draw,thick,minimum height=8mm,minimum width=10mm,align=center},
  suma/.style={draw,thick,circle,inner sep=0pt,minimum size=4.5mm}]
\node (ref) at (0,0) {$T_{\rm ref}$};
\node[suma] (s1) at (1.15,0) {};
\node[bloque] (pi) at (2.65,0) {PI\\[-1pt]{\scriptsize $K_c,\,K_i$}\\[-1pt]{\scriptsize anti-windup}};
\node[suma] (s2) at (4.05,0) {};
\node (i0) at (4.05,0.95) {$I_0$};
\node[bloque] (sat) at (5.15,0) {\begin{tikzpicture}[scale=0.22]\draw (-1.4,-1)--(-0.7,-1)--(0.7,1)--(1.4,1);\end{tikzpicture}\\[-1pt]{\scriptsize 0--2 A}};
\node[bloque] (pl) at (6.95,0) {Planta\\[-1pt]{\scriptsize Simscape}};
\node (tamb) at (6.95,1.05) {$T_{\rm amb}$};
\node (y) at (8.15,0) {$T_b$};
\draw[->] (ref) -- node[pos=0.75,below]{$-$} (s1);
\draw[->] (s1) -- node[above]{$e$} (pi);
\draw[->] (pi) -- node[above]{$\Delta I$} (s2);
\draw[->] (i0) -- (s2);
\draw[->] (s2) -- (sat);
\draw[->] (sat) -- node[above]{$I$} (pl);
\draw[->] (tamb) -- (pl);
\draw[->] (pl) -- (y);
\draw[->] (7.75,0) -- (7.75,-0.95) -- node[below]{Sensor} (1.15,-0.95) -- node[pos=0.7,right]{$+$} (s1);
\end{tikzpicture}
\caption{Lazo de control con PI, saturación y anti-windup. El error se define como $e=T_b-T_{\rm ref}$ porque la planta tiene ganancia negativa. Fuente: elaboración propia.}\label{fig:lazoPI}
\end{figure}

\subsection{Ubicación del cero}
Los polos de la planta son $-0.0477$ y $-0.0086\,\mathrm{s^{-1}}$ (constantes de tiempo de $21$ y $116\,\mathrm s$). El polo lento domina la respuesta en lazo abierto. Si el integrador se agregara sin más, el polo en el origen y el polo lento formarían un par dominante muy lento. Por eso el cero del PI se coloca sobre el polo lento, de modo que lo cancele:
\begin{equation}
z=0.008607\,\mathrm{s^{-1}}\ \Rightarrow\ T_i=\frac{1}{z}=116.2\,\mathrm s.\label{eq:cero}
\end{equation}
Con esta cancelación, la ganancia de lazo abierto se reduce a
\begin{equation}
L(s)=C(s)G_p(s)=\frac{K_c\,(0.2919\,s+0.004243)}{s\,(s+0.04772)},\label{eq:L}
\end{equation}
y la ecuación característica del lazo cerrado queda
\begin{equation}
s^2+(0.04772+0.2919K_c)\,s+0.004243K_c=0.\label{eq:carPI}
\end{equation}
Igual que en la ecuación~\eqref{eq:carP}, todos los coeficientes son positivos para $K_c>0$, así que el lazo es estable para cualquier ganancia. El polo cancelado, $-0.0086\,\mathrm{s^{-1}}$, sigue presente en el lazo cerrado. No aparece en la respuesta a la referencia, pero sí en la respuesta a perturbaciones, como se discute más adelante.

\subsection{Selección de la ganancia}
Con el cero fijo, se calculó la respuesta del lazo para $0.1\le K_c\le1.5\,\mathrm{A/K}$ en pasos de $0.05$ (Figura~\ref{fig:barrido}). El tiempo de asentamiento disminuye con la ganancia y cumple la E3 a partir de $K_c=0.65\,\mathrm{A/K}$. El esfuerzo inicial vale $\Delta I(0^+)=K_c\,|\Delta T_{\rm ref}|$, por lo que la E4 limita la ganancia a $K_c\le0.81\,\mathrm{A/K}$. Dentro del intervalo $0.65\le K_c\le0.80\,\mathrm{A/K}$ se cumplen todas las especificaciones lineales. Se escoge un valor intermedio para dejar margen hacia ambos lados:
\begin{equation}
K_c=0.75\,\mathrm{A/K},\qquad K_i=\frac{K_c}{T_i}=6.456\times10^{-3}\,\mathrm{A/(K\,s)}.\label{eq:ganancias}
\end{equation}

\begin{figure}[!tbp]\centering
\includegraphics[width=\columnwidth]{barrido_Kc}
\caption{Tiempo de asentamiento y esfuerzo de control máximo en función de $K_c$, con el cero del PI fijo en $z=0.0086\,\mathrm{s^{-1}}$. Las líneas discontinuas son las especificaciones E3 y E4. Fuente: elaboración propia con MATLAB R2025a.}\label{fig:barrido}
\end{figure}

Al sustituir $K_c=0.75\,\mathrm{A/K}$ en la ecuación~\eqref{eq:carPI} se obtienen los polos de lazo cerrado $-0.2541$ y $-0.0125\,\mathrm{s^{-1}}$, a los que se suma el polo cancelado en $-0.0086\,\mathrm{s^{-1}}$. La Figura~\ref{fig:lgrPI} los ubica sobre el lugar de las raíces. El polo en $-0.0125\,\mathrm{s^{-1}}$ queda muy cerca del cero en $-0.0145\,\mathrm{s^{-1}}$, por lo que su aporte a la respuesta es pequeño y la salida se acerca rápido a la referencia.

\begin{figure}[!tbp]\centering
\includegraphics[width=\columnwidth]{lgr_PI}
\caption{Lugar de las raíces de $K_c(s+z)G_p(s)/s$ para $0\le K_c\le2\,\mathrm{A/K}$; los cuadrados son los polos con $K_c=0.75\,\mathrm{A/K}$. Fuente: elaboración propia con MATLAB R2025a.}\label{fig:lgrPI}
\end{figure}

\subsection{Verificación en frecuencia}
La Figura~\ref{fig:bode} muestra el diagrama de Bode de $L(s)$. El margen de fase es de $98.7^\circ$ a $0.214\,\mathrm{rad/s}$ y el margen de ganancia es infinito, porque la fase nunca llega a $-180^\circ$. Se cumple la E7 con amplitud.

\begin{figure}[!tbp]\centering
\includegraphics[width=\columnwidth]{bode_PI}
\caption{Diagrama de Bode de la ganancia de lazo $L(s)=C(s)G_p(s)$ con $K_c=0.75\,\mathrm{A/K}$ y $T_i=116.2\,\mathrm s$. La línea discontinua marca la frecuencia de cruce de ganancia. Fuente: elaboración propia con MATLAB R2025a.}\label{fig:bode}
\end{figure}

\subsection{Comprobación con el modelo no lineal}
Para respaldar la elección del límite de corriente y del anti-windup se integró el modelo no lineal de la Tarea~1 con el método de Euler (paso de $0.05\,\mathrm s$) durante un arranque desde $25\,^\circ\mathrm C$ con $T_{\rm ref}=5\,^\circ\mathrm C$ (Figura~\ref{fig:arranque}). Con límite de $6.4\,\mathrm A$, la corriente se queda en el máximo y el bloque se calienta hasta $46.4\,^\circ\mathrm C$. Con límite de $2\,\mathrm A$ pero sin anti-windup, el integrador acumula error durante la saturación y el bloque baja hasta $1.92\,^\circ\mathrm C$, lo que viola la E8. Con el diseño completo, la temperatura mínima es $4.87\,^\circ\mathrm C$ y el bloque entra a la banda de $5\pm0.5\,^\circ\mathrm C$ a los $117\,\mathrm s$. La validación formal se hace con la red de Simscape en la Sección~\ref{sec:resultados}.

\begin{figure}[!tbp]\centering
\includegraphics[width=\columnwidth]{arranque_nl}
\caption{Arranque desde $25\,^\circ\mathrm C$ con el modelo no lineal para tres implementaciones del PI ($K_c=0.75\,\mathrm{A/K}$, $T_i=116.2\,\mathrm s$). Fuente: elaboración propia con MATLAB R2025a.}\label{fig:arranque}
\end{figure}

\subsection{Resumen del diseño}
La Tabla~\ref{tab:prediccion} compara el desempeño previsto del diseño con las especificaciones de la Tabla~\ref{tab:especificaciones}. Todas se cumplen. Para la implementación en el bloque \emph{PID Controller} de Simulink se usa la forma paralela con $P=K_c$, $I=K_i$ y $D=0$, saturación de salida y anti-windup por \emph{clamping} \cite{mw_pid}. Los valores se guardan en \texttt{parametros\_control.mat}.

\begin{table}[!htbp]\centering
\caption{Desempeño previsto del PI frente a las especificaciones.}\label{tab:prediccion}
{\tableformat
\begin{tabular}{clccc}\toprule
Id. & Métrica & Especificación & PI diseñado & ¿Cumple?\\\midrule
E1 & $e_{ss}$ & $\le0.1\,\mathrm K$ & $0$ & Sí\\
E2 & $M_p$ & $\le10\,\%$ & $0\,\%$ & Sí\\
E3 & $t_s$ (2\,\%) & $\le180\,\mathrm s$ & $158.6\,\mathrm s$ & Sí\\
E4 & $\Delta I_{\max}$ ($-1\,\mathrm K$) & $\le0.81\,\mathrm A$ & $0.75\,\mathrm A$ & Sí\\
E5 & Rango de $I$ & $0$ a $2\,\mathrm A$ & $0$ a $2\,\mathrm A$ & Sí\\
E6 & Pico / regreso & $\le0.5\,\mathrm K$ / $\le400\,\mathrm s$ & $0.30\,\mathrm K$ / $300\,\mathrm s$ & Sí\\
E7 & Margen de fase & $\ge45^\circ$ & $98.7^\circ$ & Sí\\
E8 & $T_{b,\min}$ / $t_{\pm0.5}$ & $\ge2\,^\circ\mathrm C$ / $\le300\,\mathrm s$ & $4.87\,^\circ\mathrm C$ / $117\,\mathrm s$ & Sí\\\bottomrule
\end{tabular}}
\par\smallskip\footnotesize Fuente: elaboración propia. E1 a E7 con el modelo lineal; E8 con el modelo no lineal. En E6 se reporta el peor caso: pico de $0.26\,\mathrm K$ por $T_{\rm amb}$ y de $0.30\,\mathrm K$ por la carga de $2\,\mathrm W$.
\end{table}
